% from ega4-16-fr.tex
\oldpage[IV-4]{41}
donc $U\mapsto\operatorname{Diff}_{U/S}^n(\mathscr{F}|U,\mathscr{G}|U)$ est
un préfaisceau de groupes additifs; en fait c'est même un \emph{faisceau}, car
pour $U$ ouvert variable dans $X$, les homomorphismes
$u\mapsto u\circ d_{U/S,\mathscr{F}|U}^n$ sont des \emph{isomorphismes} de
groupes additifs
\[
  \mathscr{H}\!om_{\mathcal{O}_U}
  (\mathscr{P}_{U/S}^n(\mathscr{F}|U),\mathscr{G}|U)
  \xrightarrow{\sim}
  \operatorname{Diff}_{U/S}^n(\mathscr{F}|U,\mathscr{G}|U),
  \tag{16.8.3.1}
  \label{IV.16.8.3.1}
\]
en vertu du fait que l'image de $\mathscr{F}$ par
$d_{X/S,\mathscr{F}}^n$ engendre $\mathscr{P}_{X/S}^n(\mathscr{F})$
(16.7.6). On note ce faisceau
$\mathscr{D}\!iff_{X/S}^n(\mathscr{F},\mathscr{G})$, et on a donc:
\end{env}

\begin{proposition}[16.8.4]
\label{IV.16.8.4}
Les isomorphismes (16.8.3.1) définissent un isomorphisme de faisceaux de
groupes additifs
\[
  \mathscr{H}\!om_{\mathcal{O}_X}
  (\mathscr{P}_{X/S}^n(\mathscr{F}),\mathscr{G})
  \xrightarrow{\sim}
  \mathscr{D}\!iff_{X/S}^n(\mathscr{F},\mathscr{G}).
  \tag{16.8.4.1}
  \label{IV.16.8.4.1}
\]
\end{proposition}

Lorsque $\mathscr{F}=\mathscr{G}=\mathcal{O}_X$, on écrit aussi
$\mathscr{D}\!iff_{X/S}^n$ au lieu de
$\mathscr{D}\!iff_{X/S}^n(\mathcal{O}_X,\mathcal{O}_X)$; il résulte de
(16.8.4) que $\mathscr{D}\!iff_{X/S}^n$ s'identifie canoniquement au
\emph{dual} du $\mathcal{O}_X$-Module $\mathscr{P}_{X/S}^n$; aussi écrit-on
$\langle t,D\rangle$ au lieu de $u(t)$ si $t$ est une section de
$\mathscr{P}_{X/S}^n$ au-dessus d'un ouvert et si $u$ est l'homomorphisme de
$\mathscr{P}_{X/S}^n$ dans $\mathcal{O}_X$ correspondant à $D$.

\begin{env}[16.8.5]
\label{IV.16.8.5}
Puisque $\mathscr{P}_{X/S}^n(\mathscr{F})$ est muni d'une structure de
$\mathcal{O}_X$-Bimodule (16.7.4), on en déduit canoniquement une structure de
$\mathcal{O}_X$-Bimodule sur
$\mathscr{H}\!om_{\mathcal{O}_X}
(\mathscr{P}_{X/S}^n(\mathscr{F}),\mathscr{G})$, donc aussi sur
$\mathscr{D}\!iff_{X/S}^n(\mathscr{F},\mathscr{G})$ en vertu de (16.8.4.1).
De façon précise, à la structure de $\mathcal{O}_X$-Module à gauche de
$\mathscr{P}_{X/S}^n(\mathscr{F})$ correspond, en vertu de la définition
(16.8.1), la structure de $\mathcal{O}_X$-Module à gauche sur
$\mathscr{D}\!iff_{X/S}^n(\mathscr{F},\mathscr{G})$ explicitée comme suit:
pour tout ouvert $U$ de $X$, toute section
$a\in\Gamma(U,\mathcal{O}_X)$ et tout opérateur différentiel
$D:\mathscr{F}|U\to\mathscr{G}|U$, $aD$ est l'opérateur différentiel qui, à
toute section $t\in\Gamma(U,\mathscr{F})$, fait correspondre la section
\[
  (aD)(t)=a(D(t))
  \tag{16.8.5.1}
  \label{IV.16.8.5.1}
\]
de $\Gamma(U,\mathscr{G})$. De même, à la structure de
$\mathcal{O}_X$-Module à droite de $\mathscr{P}_{X/S}^n(\mathscr{F})$
correspond la structure de $\mathcal{O}_X$-Module à droite sur
$\mathscr{D}\!iff_{X/S}^n(\mathscr{F},\mathscr{G})$ explicitée comme suit:
avec les mêmes notations que ci-dessus, $Da$ est l'opérateur différentiel qui,
à $t\in\Gamma(U,\mathscr{F})$, fait correspondre la section
\[
  (Da)(t)=D(at).
  \tag{16.8.5.2}
  \label{IV.16.8.5.2}
\]
\end{env}

\begin{proposition}[16.8.6]
\label{IV.16.8.6}
Si $f:X\to S$ est un morphisme localement de présentation finie,
$\mathscr{F}$ un $\mathcal{O}_X$-Module quasi-cohérent de présentation finie
et $\mathscr{G}$ un $\mathcal{O}_X$-Module quasi-cohérent, alors
$\mathscr{D}\!iff_{X/S}^n(\mathscr{F},\mathscr{G})$ est un
$\mathcal{O}_X$-Module quasi-cohérent pour l'une ou l'autre structure définies
dans (16.8.5).
\end{proposition}

La proposition résulte de ce que, sous les hypothèses faites,
$\mathscr{P}_{X/S}^n(\mathscr{F})$ est un $\mathcal{O}_X$-Module
quasi-cohérent de présentation finie (16.7.4), et de (I, 1.3.12).

\begin{env}[16.8.7]
\label{IV.16.8.7}
L'ensemble des opérateurs différentiels de $\mathscr{F}$ dans $\mathscr{G}$
(d'ordre non précisé (16.8.1)) se note
$\operatorname{Diff}_{X/S}(\mathscr{F},\mathscr{G})$; on voit encore comme
dans (16.8.3) que
$U\mapsto\operatorname{Diff}_{U/S}(\mathscr{F}|U,\mathscr{G}|U)$ est un
faisceau de groupes additifs, que nous désignerons par
$\mathscr{D}\!iff_{X/S}(\mathscr{F},\mathscr{G})$. Il est immédiat que
$\mathscr{D}\!iff_{X/S}(\mathscr{F},\mathscr{G})$ est réunion de la famille
filtrante croissante de ses sous-faisceaux
$\mathscr{D}\!iff_{X/S}^n(\mathscr{F},\mathscr{G})$; si $X$ est
quasi-compact, $\operatorname{Diff}_{X/S}(\mathscr{F},\mathscr{G})$
